% ---------- Datos de la portada ----------
\title{Actividad 2: Componentes de una computadora y procesos}
\author{Heriberto Espino Montelongo \\
  ID: 175199\\[1em]
  \text{Profesor:} Dr. Victor Giovanni Morales Murillo}
\faculty{Universidad de las Américas Puebla (UDLAP)}
\school{Licenciatura en Ciencia de Datos}
\degree{Curso: Programación Concurrente}
\date{Fecha de entrega: 20 de agosto de 2026}
\subject{Análisis de los componentes de una computadora e implementación de procesos}

\frontmatter
\maketitle
\mainmatter
\listoffigures

\begin{abbreviations}
  \item[OS] \textit{Operating System}: sistema operativo que administra los recursos de hardware.
  \item[CPU] \textit{Central Processing Unit}: procesador que controla la computadora.
  \item[RAM] \textit{Random Access Memory}: memoria principal que mantiene temporalmente instrucciones y datos.
  \item[I/O] \textit{Input/Output}: entrada y salida de datos.
  \item[PC] \textit{Program Counter}: contiene la dirección de la siguiente instrucción.
  \item[IR] \textit{Instruction Register}: contiene la instrucción que está siendo ejecutada.
  \item[MAR] \textit{Memory Address Register}: especifica la dirección de memoria que será leída o escrita.
  \item[MBR] \textit{Memory Buffer Register}: contiene los datos que se escriben o leen de la memoria.
  \item[I/O AR] \textit{Input/Output Address Register}: especifica el dispositivo de I/O.
  \item[I/O BR] \textit{Input/Output Buffer Register}: intercambia datos con un I/O module.
  \item[AC] \textit{Accumulator}: register utilizado como almacenamiento temporal.
  \item[CISC] \textit{Complex Instruction Set Computer}: arquitectura de procesador con un conjunto de instrucciones complejas e integradas.
  \item[RISC] \textit{Reduced Instruction Set Computer}: arquitectura de procesador simplificada orientada a la ejecución rápida de instrucciones básicas.
  \item[x86] Arquitectura de juego de instrucciones CISC desarrollada originalmente por Intel para procesadores de 32 y 64 bits (x86\_64).
  \item[ARM] \textit{Advanced RISC Machine}: arquitectura de procesamiento basada en diseño RISC de alto rendimiento y bajo consumo energético.
  \item[PID] \textit{Process Identifier}: identificador único asignado por el kernel a un proceso en ejecución.
  \item[PPID] \textit{Parent Process Identifier}: identificador de proceso correspondiente al proceso padre que generó la ejecución actual.
  \item[UID] \textit{User Identifier}: número entero que identifica de forma persistente al usuario propietario del proceso.
  \item[EUID] \textit{Effective User Identifier}: identificador de usuario efectivo que el sistema evalúa para otorgar permisos y accesos en tiempo de ejecución.
\end{abbreviations}
